% Strategic Management – Midterm Study Guide — EU business casual
% Compile with: pdflatex midterm.tex

\documentclass[8pt,landscape,a4paper]{article}
\usepackage[utf8]{inputenc}
\usepackage[T1]{fontenc}
\usepackage[left=0.8cm,right=2.6cm,top=0.8cm,bottom=0.8cm,marginparwidth=2cm,marginparsep=0.3em]{geometry}
\usepackage{multicol}
\usepackage{marginnote}
\usepackage{enumitem}
\setlist{nosep,leftmargin=*}
\usepackage{amsmath}
\usepackage{booktabs}
\usepackage{parskip}
\usepackage{microtype}
\usepackage{xcolor}
\usepackage{fancyhdr}
\definecolor{eublue}{RGB}{0,51,153}
\definecolor{charcoal}{RGB}{51,51,51}
\definecolor{softgrey}{RGB}{240,240,240}
\definecolor{notegrey}{RGB}{90,90,90}
\newcommand{\handnote}[1]{\marginnote{\raggedright\footnotesize\color{notegrey}\textit{#1}\vspace{0.4em}}}
\newcommand{\hlpink}[1]{\textcolor{charcoal}{\textbf{#1}}}
\newcommand{\hlpeach}[1]{\textcolor{eublue}{\textbf{#1}}}
\setlength{\parskip}{0.2em}
\setlength{\columnsep}{0.5em}
\pagestyle{fancy}
\fancyhf{}
\fancyfoot[L]{\raisebox{0.8em}{\small\color{notegrey} Strategic Mgmt}}
\fancyfoot[R]{\raisebox{0.8em}{\small\color{notegrey}\thepage}}
\fancypagestyle{plain}{\fancyhf{}\fancyfoot[L]{\raisebox{0.8em}{\small\color{notegrey} Strategic Mgmt}}\fancyfoot[R]{\raisebox{0.8em}{\small\color{notegrey}\thepage}}}
\raggedright

\begin{document}
\begin{center}
\vspace{0.2em}
{\normalsize\bfseries\color{eublue} Strategic Management -- Midterm Study Guide}\\[0.2em]
\footnotesize\textcolor{notegrey}{Use for quick recall; in exam: define terms, give one example, tie to competitive advantage.}\\[0.35em]
\noindent\colorbox{softgrey}{\parbox{0.92\textwidth}{\centering\footnotesize\color{charcoal}
\textbf{Key:} \textbf{Definitions} \quad \textcolor{eublue}{\textbf{Frameworks \& tips}} \quad \textit{Margin notes}
\vspace{0.25em}}}
\end{center}
\vspace{-0.4em}

\begin{multicols}{2}

\section*{\color{eublue}Unit 0 -- Course \& Big Ideas}\handnote{Long-term, scarce resources -- know these by heart.}
\begin{itemize}[leftmargin=*,nosep]
\item \textbf{\hlpink{Strategy}}: long-term, hard-to-reverse decisions; compete for \textbf{\hlpink{scarce resources}} (customers, talent, capital). Limited ``spots'' in the market; we compete for them.
\item \textbf{Strategic vs tactical}: Strategic = long-term, shapes future, costly to undo (what to study, where to live, career change, mergers, entering a country). ``No second chance'' for many. Tactical = short-term, day-to-day (scheduling, small price changes).
\item \textbf{Forecasting}: future uncertain (economy, politics, climate, health); we imagine \textbf{scenarios}, we don't know exact outcomes. Theory is a \textbf{tool}, not the answer.
\item \textbf{Key rules}: \textbf{\hlpink{Know yourself}} = strengths, weaknesses, goals, constraints. \textbf{\hlpink{Know your environment}} = opportunities and threats (macro + industry).
\item \textbf{Social science}: politics, power, culture, stakeholders matter. Rational analysis has limits. No single ``right'' answer.\handnote{Not just shareholders! Workers, society, env.}
\item \textbf{Stakeholders} (not just shareholders): workers, community, society, environment, state/family owners; long-term success needs more than profit. Debate: \textbf{shareholder primacy} vs \textbf{stakeholder/CSR/sustainability}.
\end{itemize}

\section*{\color{eublue}Unit 1 -- Strategy \& The Art of War}\handnote{Apply: terrain? enemy? what does firm know?}
\begin{itemize}[leftmargin=*,nosep]
\item \textbf{Strategos} (Greek) = the general; art of commanding the army. Linked to war \& competition for limited positions and resources.
\item \textbf{Why ``war''?} On a battlefield, armies compete for territory; in business, firms compete for customers, talent, capital. Same logic: limited spots, someone wins, someone loses.
\item \textbf{Strategic decisions in life}: what to study, career, where to live, partner, children $\Rightarrow$ big effects, no easy second chance. Think 10+ years ahead.
\item \textbf{Sun Tzu analogies}: Battlefield $\to$ \textbf{market} (where competition happens). Enemy $\to$ \textbf{competitors/obstacles}. General $\to$ \textbf{leader/strategist}. Terrain $\to$ \textbf{industry structure, rules, geography}.
\item \textbf{Core lesson}: If you \hlpeach{know yourself} (resources, capabilities, goals), \hlpeach{know the enemy} (competitors, their moves), and \hlpeach{know the terrain} (market, industry, regulation), your chances of success are much higher.
\item \textit{\hlpeach{Exam tip}}: In a case, ask: What is the terrain? Who are the enemies? What does the firm know or not know?
\end{itemize}

\section*{\color{eublue}Unit 2 -- Objectives, Values, Performance, SWOT}\handnote{Mission = who. Vision = where. Objectives = how.}

\subsection*{A. Strategy as intention}
Strategy = where we want to go + how we get there. Long-term = often $>$10 years.\handnote{Always compare performance to something!}

\subsection*{B. Mission, vision, objectives}
\begin{itemize}[leftmargin=*,nosep]
\item \textbf{\hlpink{Mission}} = Who am I? What do I do? \textbf{Current state.} Products/services + target market(s). One sentence (e.g.\ ``I produce and sell cars'').
\item \textbf{\hlpink{Vision}} = What do I want to become? \textbf{Future state}, long-term (e.g.\ ``Leader in electric vehicles in Europe by 2035'').
\item \textbf{\hlpink{Strategic objectives}} = Concrete steps from mission to vision (milestones, targets).
\item \textbf{Objectives} (general): profit, growth, survival, reputation. \textbf{Values}: what firm will/won't do (ethics, CSR). Constrain \textit{how} you compete.
\item \textbf{Performance}: always ``compared to what?'' \textbf{Referent group} = similar firms in industry. \textbf{History} = firm's own past. \textbf{Intention} = chosen constraints (e.g.\ closing Sundays).
\end{itemize}

\subsection*{C. Measuring performance}
\begin{itemize}[leftmargin=*,nosep]
\item \textbf{Accounting} (non-listed/family): \textbf{\hlpink{ROA}} = profit/total assets (efficiency); \textbf{\hlpink{ROE}} = profit/equity (return for shareholders). A number only makes sense vs industry average, main rival, or own past. No absolute ``good'' number.
\item \textbf{Market} (listed): Shareholder return = (price end $-$ price start) + dividends. \textbf{Opportunity cost} = next best alternative (other firms, index, risk-free). Satisfied when return $\ge$ opportunity cost.
\item \textbf{\hlpink{Cost of equity}} = minimum return shareholders require: $\text{Cost of Equity} = R_f + \beta \times (\text{Market Return} - R_f)$. \textbf{$R_f$} = risk-free rate (e.g.\ govt bonds). \textbf{$\beta$} = sensitivity to market. \textbf{Market Return} = expected market return (e.g.\ index). \hlpink{Satisfaction}: return $>$ cost of equity (or better than industry average).
\end{itemize}\handnote{Formula: Rf + beta (Market - Rf). Compare vs industry.}

\subsection*{D--E. Competitive advantage \& Porter}
\begin{itemize}[leftmargin=*,nosep]
\item \textbf{\hlpink{CA}} at each \textbf{transaction} when customer chooses you instead of a rival or not buying. Built sale by sale; over many transactions it looks ``sustained''.
\item \textbf{Low-cost}: win on price; customers price-sensitive (elastic). E.g.\ budget airlines, generic products.
\item \textbf{Differentiation}: customers strongly prefer your offer (brand, quality, design, service); inelastic. E.g.\ luxury brands, Apple.
\item \textbf{Focus}: narrow segment (niche); can combine with low-cost or differentiation. \textbf{Stuck in the middle}: Porter warned---neither low-cost nor differentiated $\Rightarrow$ worst of both (low margins, no loyalty).
\end{itemize}\handnote{Low-cost, diff, focus. Never stuck in the middle!}

\subsection*{F. Governance \& CSR}
Companies = social organisations. Stakeholders: shareholders (return), managers (power, pay), workers (wages, conditions), customers (quality, price), society (jobs, taxes, pollution). \textbf{Agency problem}: managers $\neq$ owners; need carrots (stock options, bonuses) and sticks (board oversight). \textbf{Externalities}: negative (e.g.\ pollution) and positive (jobs, taxes) often not fully priced. Debate: shareholder primacy vs stakeholder/CSR/sustainability (climate, social expectations).\handnote{Agency: managers vs owners. Externalities = not priced.}

\subsection*{G. SWOT done right}
S/W = \textbf{internal} (resources, capabilities). O/T = \textbf{external} (environment, industry, PESTEL). \textbf{Bad}: vague (``our people are our strength'', ``the market is growing'')---no link to why customers choose you. \textbf{Good}: specific strengths/weaknesses tied to CA (e.g.\ ``logistics allows 24h delivery'', ``brand in segment X'', ``location next to university''); concrete O/T (e.g.\ ``new regulation favours renewables'', ``two new entrants 2024''). \textit{\hlpeach{Exam tip}}: What do we have (S/W)? What's outside (O/T)? How do S and O combine?\handnote{S/W = internal. O/T = external. Be specific!}

\section*{\color{eublue}Unit 3 -- PESTEL, Five Forces, Clusters}\handnote{Macro + industry. Can't control but must adapt.}

\subsection*{Environment}
External factors firm \textbf{cannot control} but that \textbf{influence} behaviour and results. Managers' job: detect opportunities and threats. Two levels: \textbf{macro} (country/world) and \textbf{industry} (specific arena).

\subsection*{\hlpeach{PESTEL}}
\begin{itemize}[leftmargin=*,nosep]
\item \textbf{P}olitical: govt, stability, policy, taxes, trade, war/peace. E.g.\ change of government $\to$ new subsidies or tariffs.
\item \textbf{E}conomic: growth, inflation, interest rates, unemployment, exchange rates. E.g.\ recession $\to$ less demand; strong euro $\to$ export harder.
\item \textbf{S}ocial: demographics, culture, lifestyles, education, inequality. E.g.\ ageing population $\to$ health demand; remote work $\to$ tech demand.
\item \textbf{T}echnological: R\&D, new tech, digitalisation. E.g.\ AI, automation $\to$ new products, displaced jobs.
\item \textbf{E}cological: climate, resources, pollution. E.g.\ regulation on emissions $\to$ cost or green tech opportunities.
\item \textbf{L}egal: labour, consumer, competition, environment, IP. E.g.\ GDPR, minimum wage.
\end{itemize}
PESTEL changes $\Rightarrow$ new opportunities \& threats $\Rightarrow$ firms must adapt.\handnote{Six letters: P E S T E L. One example each.}

\subsection*{\hlpeach{Porter's Five Forces}}
Aim: \textbf{industry attractiveness} (profit potential). Strong forces = less attractive; weak = more attractive. Always from \textbf{firm's perspective} (good for buyers/suppliers = bad for firm).
\begin{enumerate}[leftmargin=*,nosep]
\item \textbf{Rivalry}: high when many competitors, low growth, commoditised (low differentiation), high exit barriers, high fixed costs. Result: price wars, lower margins.
\item \textbf{New entrants}: easy entry $\Rightarrow$ more firms $\Rightarrow$ more rivalry. \textbf{Barriers}: natural (economies of scale, know-how, tech, learning curve); artificial (licenses, regulation, tariffs, patents); capital requirements; access to distribution; strong brands; expected retaliation.
\item \textbf{Substitutes}: products/tech that satisfy same need (e.g.\ train vs plane; streaming vs cinema). Disruptive technology or consumer behaviour can increase this force.
\item \textbf{Buyer power}: high when few large buyers (concentration); low switching cost; they have full info on your costs; they can backward integrate (make it themselves). Result: they push down price $\Rightarrow$ less attractive for you.
\item \textbf{Supplier power}: high when few suppliers; you depend on their specific input (differentiation); they have info on your business; they can forward integrate (become your competitor). Result: they push up input prices $\Rightarrow$ less attractive.
\end{enumerate}
Goal: \textbf{\hlpeach{no organising power}} for either buyers or suppliers (keep margin and control).\handnote{Strong force = bad for you. Weak = good.}

\subsection*{Limits \& extras}
Performance = \textbf{\hlpeach{industry effect}} (attractiveness) + \textbf{\hlpeach{firm effect}} (own strategy and resources). Firm can do well in unattractive industry (strong resources) or badly in attractive one (weak strategy). Not all five forces equally relevant; focus on strongest 2--3. \textbf{Complements}: products that go with yours (e.g.\ software + hardware) can increase demand. \textbf{Sixth force}: public admin/regulators (Europe: energy, telecoms, health). \textbf{Segments \& strategic groups}: not all players face same forces. Five Forces = simplification of reality.\handnote{Industry + firm effect. Not all forces equal -- pick 2--3.}

\subsection*{Clusters}\handnote{Geography! Firms + suppliers + unis. Silicon Valley, Toyota city.}
\textbf{Geographical agglomeration} of firms in one activity + suppliers + related industries + universities/research + support (e.g.\ Silicon Valley, Toyota city, Italian textile districts). Benefits: talent attraction, entrepreneurship, technological innovation, access to finance and networks. Proximity reduces transaction costs and speeds learning.

\subsection*{Industry definition \& attractiveness}
Define from \textbf{customer's point of view}: not just ``all car makers'' but all offers that \textbf{meet the same need/wish} (including substitutes, luxury experiences). \textbf{Attractive} when opportunities $>$ threats $\Rightarrow$ easier profit, growth, survival. \textbf{Unattractive} when threats $>$ opportunities. Manager's job: enter attractive industries, avoid/exit bad ones when possible.\handnote{Customer need = industry. Attractive = O $>$ T.}

\subsection*{Climate (PESTEL case)}\handnote{Arctic example: all 6 PESTEL.}
Climate change = economic and strategic, not only environmental. Melting Arctic/Greenland: new sea routes, trade patterns; resources (oil, gas, rare earths) more accessible; geopolitical competition, legal disputes. PESTEL: Political (sovereignty, military), Economic (routes, costs), Social (local communities), Technological (ice-breaking, extraction), Environmental (fragile ecosystem), Legal (UNCLOS, maritime law).

\section*{\color{eublue}Unit 4 -- \hlpeach{RBV} \& Internal Resources}\handnote{Internal view! 70-80\% = you, not industry.}
\begin{itemize}[leftmargin=*,nosep]
\item \textbf{RBV} (Barney and others): strategy from \textbf{what firm has inside} (resources \& capabilities), not only external (PESTEL, Five Forces). Internal view \textbf{complements} external view.
\item \hlpeach{70--80\%} of success/failure = \textbf{internal} factors (management, people, resources, capabilities), not just industry or macro.
\item You don't need to be ``the best''; need to be \textbf{better than competitors} and \textbf{sustain} that advantage over time. \textbf{\hlpink{Sustainable CA}} = advantage that lasts because rivals cannot easily copy or substitute it.
\item \textbf{Resources} = what you \textit{have} (assets, people, brand, patents). \textbf{Capabilities} = what you \textit{do} with them (processes, routines, coordination). Both matter for RBV.
\end{itemize}\handnote{Have = resources. Do = capabilities.}

\subsection*{\hlpeach{Value Chain} (Porter, 1985)}\handnote{Where value is created.}
Breakdown of firm into activities. \textbf{Primary}: inbound logistics, operations, outbound logistics, marketing, service. \textbf{Supporting}: procurement, HR, technology, infrastructure. \textbf{Value system} = firm + suppliers + customers. Advantage from: (1) activities themselves; (2) \textbf{horizontal links} (coordination, optimization inside firm); (3) \textbf{vertical links} (with suppliers/buyers). Use: identify S/W; outsourcing vs core competence; sources of SCA. E.g.\ Mercadona, Ford Valencia.

\subsection*{\hlpeach{VRIN / VRIO}}\handnote{All 4 + Org needed for sustained CA.}
\textbf{V}aluable $\Rightarrow$ helps exploit opportunities or neutralise threats (creates value). \textbf{R}are $\Rightarrow$ few or no competitors have it. \textbf{I}nimitable $\Rightarrow$ difficult or costly to copy (e.g.\ unique culture, tacit knowledge). \textbf{N}on-substitutable $\Rightarrow$ cannot be replaced by something else that does the same job. \textbf{O} (VRIO) = \textbf{Organisation}: structure and processes to \textbf{exploit} the resource (otherwise V+R+I+N wasted). Only resources that meet all four (V+R+I+N) and are well organised (O) support \textbf{sustainable} CA. If valuable but not rare, or rare but easy to copy, advantage is temporary.

\subsection*{Tangible vs.\ intangible}
\begin{center}
\footnotesize
\begin{tabular}{@{}lll@{}}
\toprule
\textbf{Aspect} & \textbf{Tangible} & \textbf{Intangible} \\
\midrule
What & Physical, financial & Human, tech, reputation, brand \\
Where & In financial statements & Often not in books \\
Control & Easier & Harder; ``double loss'' when people leave \\
\bottomrule
\end{tabular}
\end{center}
\hlpeach{$\sim$70\%} of value can be intangibles.\handnote{Intangibles often not in books!}

\subsection*{Human-based resources}
Knowledge, motivation, behaviour, loyalty ``in the person''; firm does \textbf{not own} the person. They can leave. When key people leave $\Rightarrow$ \textbf{\hlpink{double loss}}\handnote{Person + knowledge gone!} (you lose both the person and the knowledge; what was in their heads is gone). \textbf{Retention}: culture, identity, belonging, incentives (e.g.\ ``we are champions''), not only pay. If resource is \textbf{complementary} to others and hard to imitate, competitors cannot easily ``buy'' the same advantage (e.g.\ star player + team + culture + system). \textbf{Example---football}: Real Madrid, Barça, Man City; players = human resources; Ronaldo, Messi = rare, VRIO. Reputation (club brand) goes up or down with results and behaviour. Most value \textbf{intangible} (not on balance sheet).

\subsection*{Reputation / brand}\handnote{Brand up or down with behaviour.}
Use of brand name can \textbf{increase or decrease} value over time (quality, scandals, consistency). Intangible; not fully reflected in financial statements. Same for human capital and many capabilities.

\section*{\color{eublue}Case 1 -- Madonna \& Lola Flores}\handnote{Economic vs social success. Both = human resources.}

\subsection*{Economic vs social success}
\textbf{Madonna}: one of the most profitable artists; \textbf{economic} success (revenue, profitability, global reach). \textbf{Lola Flores}: huge \textbf{social and cultural} contribution---working class, radio, fashion, identity; ``did more than profit'' for society. Both = \hlpink{human-based strategic resources}; we can assess value and rareness in economic \& social terms.

\subsection*{\hlpeach{Grant's four elements} (apply to any case)}\handnote{Goals, env, resources VRIO?, implementation.}
(1) \textbf{Simple, consistent, long-term goals}. (2) \textbf{Profound understanding of the environment} (industry, culture). (3) \textbf{Objective appraisal of resources and capabilities}. (4) \textbf{Effective implementation} (organisation, choices, actions).

\subsection*{Madonna}
\textbf{Goals}: ``Queen of Pop'', global cultural icon. \textbf{Environment}: New York; talent not enough; industry intensely competitive, volatile; need differentiation and reinvention. \textbf{Resources}: discipline, charisma, work ethic; reinvention of personas/styles; strong personal brand; network (choreographers, designers, producers); instinct for trends. \textbf{Implementation}: autonomy over artistic direction; controversy \& media to stay visible; teams for each phase. \textbf{VRIO examples}: reinvention ability (valuable, rare, hard to copy); global brand; creative network; instinct for trends + discipline.\handnote{Reinvention = VRIO. Differentiation.}

\subsection*{Lola Flores}
\textbf{Goals}: succeed from modest beginnings; central Spanish cultural figure. \textbf{Environment}: economic hardship, limited roles for women; differentiated by blending flamenco, copla, popular theatre. \textbf{Resources}: charisma, stage presence; ``La Faraona'' identity (gipsy roots); hybrid style; family dynasty (``Flores Clan''). \textbf{Implementation}: film, radio, TV; media presence, humour, authenticity; family in performances. \textbf{VRIO examples}: charisma \& emotional connection; ``La Faraona'' identity; fusion style (flamenco + copla + theatre); family clan amplifying brand.\handnote{Identity + family = rare, hard to copy.}

\subsection*{\hlpeach{Exam tip}}\handnote{Goals $\to$ env $\to$ VRIO $\to$ implementation.}
Apply Grant's four to any case: (1) What are the goals? (2) How well does the firm understand the environment? (3) What resources/capabilities (VRIO)? (4) How effective is implementation?

\section*{\color{eublue}Quick revision -- Can you explain\ldots?}\handnote{Run through these before the exam!}
\begin{itemize}[leftmargin=*,nosep]
\item \textbf{Strategy} in one sentence (long-term, scarce resources, competition).
\item \textbf{Strategic vs tactical} with one example each.
\item \textbf{Art of War}: know yourself, enemy, terrain $\Rightarrow$ apply to a firm.
\item \textbf{Mission vs vision vs objectives} (current state, future state, steps).
\item \textbf{ROA vs ROE} and when to use accounting vs market valuation.
\item \textbf{Cost of equity} formula and what ``shareholders satisfied'' means.
\item \textbf{CA} at the transaction level (customer chooses you).
\item \textbf{Porter's three strategies} (low-cost, differentiation, focus) with examples.
\item \textbf{SWOT done right} (specific, linked to why customers choose you).
\item \textbf{PESTEL}---six factors and one example each.
\item \textbf{Five Forces}---all five; what makes each force strong (bad for you)?
\item \textbf{Value chain}: primary + supporting; horizontal + vertical links.
\item \textbf{Barriers to entry}: natural vs artificial.
\item \textbf{Buyer and supplier power}: concentration, differentiation, information, vertical integration.
\item \textbf{Industry attractiveness} and ``industry effect + firm effect''.
\item \textbf{RBV}: why internal resources matter (70--80\%); \textbf{VRIN/VRIO}.
\item \textbf{Tangible vs intangible} (what, where in statements, control).
\item \textbf{Human resources}: why ``double loss'' when they leave; retention.
\item \textbf{Madonna vs Lola Flores}: economic vs social success; both as VRIO human resources.
\end{itemize}\handnote{If you can explain all these, you're ready.}

\section*{\hlpeach{Common question angles}}\handnote{Memorise these 4 question types!}
\begin{itemize}[leftmargin=*,nosep]
\item \textbf{``Compare internal and external analysis.''} External = PESTEL + Five Forces (opportunities, threats, industry attractiveness). Internal = RBV, resources, capabilities, VRIO (strengths, weaknesses). Both needed for strategy.
\item \textbf{``Why might a firm have CA in an unattractive industry?''} Strong \textbf{internal} resources and capabilities (VRIO) can outweigh a bad industry. E.g.\ niche player with unique know-how.
\item \textbf{``Why do financial statements underestimate firm value?''} Many key resources are \textbf{intangible} (human capital, reputation, culture) and not fully reflected; $\sim$70\% of value can sit in intangibles.
\item \textbf{``Apply Grant's four elements to [case].''} Goals $\to$ environment understanding $\to$ resources/capabilities (VRIO?) $\to$ implementation.
\end{itemize}

\section*{\color{eublue}Exam-style patterns and example answers}

\subsection*{\hlpeach{General sentence formula}}
\begin{itemize}[leftmargin=*,nosep]
\item \textbf{Pattern:} \emph{[Concept / force] is [high/low] because [reason], which makes the industry/situation [attractive/unattractive] for [who], therefore [strategic implication].}
\item \textbf{Name the framework first:} ``From a Porter Five Forces perspective\ldots'', ``Using RBV/VRIO\ldots'', ``At the macro level (PESTEL)\ldots''.
\item Use \textbf{economic language}: profit potential, industry structure, bargaining power, switching costs, opportunity cost, non-recoverable cost, double loss, agency problem, industry effect vs firm effect.
\item Be \textbf{specific}, not generic: concrete examples (rail as substitute, Airbus/Boeing duopoly, certification costs), not just ``politics/economy/tech''.
\item Always finish with \textbf{``therefore\ldots''}: link analysis to a conclusion (attractive/unattractive, CA/no CA, what the firm/industry should do).
\item If possible, \textbf{mirror the question angle}: e.g.\ ``This question is essentially about comparing internal and external analysis\ldots'' or ``Even in an unattractive industry, a firm can sustain CA if it has VRIO resources\ldots''.
\end{itemize}

\subsection*{\hlpeach{Reluctance to pay for training (human capital / RBV)}}
Employees are \textbf{human-based strategic resources}, but the firm does not own the person, only their labour. When a company finances education or certification, it incurs a non-recoverable cost. If the employee later leaves, the firm suffers a \textbf{double loss}: it loses both the worker and the additional knowledge it paid to develop. This risk is highest for \textbf{general} (transferable) qualifications that increase the employee's market value for rivals. Firms therefore tend to underinvest in general training and favour firm-specific skills or strong retention mechanisms (culture, identity, incentives).

\subsection*{\hlpeach{EU airline industry attractiveness (Five Forces + PESTEL)}}
From a Five Forces/PESTEL view, the EU airline industry has \textbf{low profit potential}:
\begin{itemize}[leftmargin=*,nosep]
\item \textbf{Supplier power (high):} aircraft supply is an effective duopoly with Airbus/Boeing; high switching costs (fleet standardisation, certifications, long-term contracts) entrench dependence.
\item \textbf{Substitutes (high):} high-speed rail satisfies the same need on short/medium routes; puts downward pressure on fares.
\item \textbf{Rivalry (high):} high fixed costs and low differentiation generate intense price competition.
\item \textbf{Buyer power (high):} aggregators create full price transparency and low switching costs.
\end{itemize}
Overall, threats outweigh opportunities: structurally, the industry is \textbf{unattractive}.

\subsection*{\hlpeach{Making an unattractive industry more attractive}}
Industry attractiveness improves when the strongest forces are weakened or better managed:
\begin{itemize}[leftmargin=*,nosep]
\item \textbf{Reduce rivalry:} consolidation (mergers, alliances) lowers the number of effective competitors and softens price wars.
\item \textbf{Mitigate supplier power:} long-term contracts and diversified sourcing partly rebalance bargaining power with aircraft suppliers.
\item \textbf{Differentiate:} loyalty programmes, brand and service quality create more price-inelastic demand, reducing the impact of substitutes and buyer power.
\end{itemize}
The logic is always the same: \textbf{identify the strongest forces, explain why they are high, and propose strategies that weaken them.}

\subsection*{\hlpeach{Extra phrases to drop into answers}}
\textbf{Internal vs external analysis.} At the external level, PESTEL and Five Forces describe the \textbf{structure and profit potential of the industry}; at the internal level, RBV and VRIO explain why, given the same environment, some firms systematically outperform others. A convincing diagnosis must combine both views: external analysis to understand the playing field, and internal analysis to see whether the firm has the means to build and sustain competitive advantage on that field.

\textbf{CA in an unattractive industry.} An industry can be structurally unattractive (strong rivalry, powerful suppliers and substitutes), yet a firm may still earn above-average returns if it controls VRIO resources (e.g.\ brand, technology, network effects) that competitors cannot easily replicate. In that case, the negative industry effect is more than offset by a positive firm effect; exit is not always necessary if a protected niche exists.

\textbf{Intangibles \& financial statements.} Conventional financial statements systematically understate firm value because they record tangible assets at historical cost while most of the competitive edge resides in intangibles such as human capital, reputation, culture and routines. These assets are rarely fully capitalised, yet they are decisive for VRIO; the ``missing'' value is embedded in resources the accounts do not see.

\textbf{Grant's four elements (generic template).} In any case: (1) identify goals (long-term economic, social, reputational objectives); (2) assess understanding of the environment (industry structure, macro trends, culture); (3) evaluate resources and capabilities under VRIO; (4) examine implementation (are choices coherent with goals, environment and resource base?). Where these four elements are aligned, strategy is coherent and sustainable; where they diverge, performance suffers.

\textbf{Airlines: structural context.} Given capital intensity, regulation and exposure to powerful suppliers and substitutes, airlines operate in a high-risk, low-margin strategic context, which forces them to compete on efficiency and search for pockets of differentiation (network design, hubs, brand, loyalty) rather than on input prices alone.

\end{multicols}

\newpage
\vspace*{0.3em}
\begin{center}
\textbf{\color{eublue}Quick recap -- last page}\\[0.3em]
\small\textcolor{notegrey}{(Section summaries, not margin notes)}
\end{center}
\vspace{0.4em}
\noindent\color{notegrey}
\setlength{\parskip}{0.3em}
\begin{multicols}{2}

\textbf{Unit 0 -- Course \& Big Ideas}\\
Long-term, scarce resources. Not just shareholders: workers, society, env.

\textbf{Unit 1 -- Strategy \& The Art of War}\\
Apply: terrain? enemy? what does firm know?

\textbf{Unit 2 -- Objectives, Performance, SWOT}\\
Mission = who. Vision = where. Objectives = how. Formula: Rf + beta (Market - Rf). Low-cost, diff, focus. S/W = internal. O/T = external. Be specific!

\textbf{Unit 3 -- PESTEL, Five Forces, Clusters}\\
Macro + industry. P E S T E L. Strong force = bad. Industry + firm effect. Clusters: geography. Customer need = industry. Attractive = O $>$ T.

\textbf{Unit 4 -- RBV \& Internal Resources}\\
70-80\% = you. Value chain: primary + supporting; horizontal + vertical links. Have = resources. Do = capabilities. VRIO + Org. Intangibles not in books. Double loss when people leave. Brand up or down.

\textbf{Case -- Madonna \& Lola}\\
Both = human resources. Grant's four: goals $\to$ env $\to$ VRIO $\to$ implementation. Reinvention vs identity+family.

\textbf{Quick revision}\\
Run through before exam. If you can explain all these, you're ready.

\textbf{Common question angles}\\
Compare internal/external. CA in unattractive industry? Intangibles. Grant's four.

\end{multicols}
\vfill
\normalsize\normalfont\color{black}
\end{document}
